\documentclass[11pt]{article}
\usepackage[a4paper,margin=25mm]{geometry}
\usepackage{amsmath,amssymb,hyperref}
\title{Incompatible density claims in conjecture 00000008487}
\author{ziangni-sys}
\date{4 October 2026}
\begin{document}
\maketitle
Both versions assert that the locked-function set is dense $G_\delta$
and that its complement is dense and open. These assertions cannot
hold simultaneously in any nonempty topological space. The contradiction
does not depend on the locking convention, measure or fractal dimension.

\paragraph{Topological obstruction.}
Let $Y$ be a nonempty topological space and $L\subseteq Y$ be dense.
If $Y\setminus L$ is open and nonempty, density of $L$ implies
\[
L\cap(Y\setminus L)\ne\varnothing,
\]
which is impossible. Hence an open complement of a dense set must be
empty. But an empty set cannot be dense in a nonempty space: its closure
is empty. We have therefore proved
\[
\neg\bigl(L\text{ dense}\ \land\ Y\setminus L\text{ open}
\ \land\ Y\setminus L\text{ dense}\bigr).
\]
This argument is valid without a Baire-space assumption, and even
without the additional $G_\delta$ assertion.

\paragraph{Application to the actual interval function space.}
Take the genuine interval $I=[0,1]\subset\mathbb R$ and its continuous
observable space $Y=C(I,\mathbb R)$, with its usual uniform topology.
This space is nonempty, since the constant zero function belongs to it.
Fix any interval map, for example the continuous identity map
$T(x)=x$, and let $L$ be its actual set of locked observables according
to the intended definition. Since the obstruction holds for \emph{every}
subset of $Y$, it holds for this actual $L$. Thus $L$ cannot satisfy the
two asserted density clauses simultaneously. This disproves the
conjecture as written, without imposing a definition of locking that
might alter the statement.

A dense $G_\delta$ set can certainly have a dense complement; the added
\emph{openness} of that complement is what causes the contradiction.
No conclusion about the separate Lebesgue-measure or box-dimension
claims is needed. The proof remains valid for other nonempty spaces
of observables and any topology on them.

\paragraph{Lean correspondence.}
The accompanying project proves the obstruction for every actual
nonempty topological space and every subset, using Mathlib's
\texttt{Dense} and \texttt{IsOpen} predicates. It specializes to
\texttt{ContinuousMap (Set.Icc (0:Real) 1) Real}, supplies the actual
zero observable, and proves that no subset can satisfy the complete
claimed topology conjunction, including the $G_\delta$ predicate. This universal
statement covers the actual locked-function set without an assumed
certificate or a replacement dynamics model. The project is pinned
to Lean 4.19.0 and Mathlib commit
\texttt{c44e0c8ee63ca166450922a373c7409c5d26b00b}.
\end{document}
